% Transport parallele : deux vecteurs tangents deplaces le long d'une
% geodesique, l'angle qu'ils forment etant conserve.
% Utilise dans la \cref{sec:context-riemann-notion}.
% Styles et briques communs aux figures de geometrie riemannienne du
% \cref{ch:context-riemann}. Ce fichier peut etre \input plusieurs fois.
\definecolor{rmanifold}{RGB}{90,90,90}
\definecolor{rtangent}{RGB}{50,110,120}
\definecolor{rgeodesic}{RGB}{40,70,140}
\definecolor{rvector}{RGB}{175,95,45}
\tikzset{
  riemann/.style={
    >={Stealth[length=4.5pt,width=3.5pt]},
    line join=round,
    font=\footnotesize,
  },
  % La nappe qui represente la variete, centree sur l'origine.
  rsurface/.pic={
    \fill[rmanifold!7]
      (-3.2,0.15) .. controls (-2.95,1.75) and (-1.5,2.35) .. (0,2.35)
      .. controls (1.5,2.35) and (2.95,1.7) .. (3.2,0.15)
      .. controls (1.9,-0.7) and (-1.9,-0.7) .. (-3.2,0.15) -- cycle;
    \draw[rmanifold!70, thick]
      (-3.2,0.15) .. controls (-2.95,1.75) and (-1.5,2.35) .. (0,2.35)
      .. controls (1.5,2.35) and (2.95,1.7) .. (3.2,0.15);
    \draw[rmanifold!70, thick]
      (-3.2,0.15) .. controls (-1.9,-0.7) and (1.9,-0.7) .. (3.2,0.15);
  },
  % Le plan tangent, a placer en un point donne via un shift.
  rplane/.pic={
    \fill[rtangent!10, draw=rtangent!65]
      (-1.35,-0.42) -- (1.55,-0.88) -- (2.05,0.22) -- (-0.85,0.68) -- cycle;
  },
}

\begin{tikzpicture}[riemann]
  \coordinate (P) at (-1.85,1.05);
  \coordinate (Q) at ( 1.65,1.50);

  \path (0,0) pic {rsurface};
  \node[rmanifold!85] at (2.75,0.75) {$\Manifold$};

  % la geodesique support
  \draw[rgeodesic, very thick]
    (P) .. controls (-0.75,1.90) and (0.60,2.05) .. (Q);
  \node[rgeodesic!85, anchor=north, inner sep=3pt] at (-0.15,1.55)
    {$\transport_{\param_0\to\param_1}$};

  % le couple de vecteurs en theta_0
  \draw[rvector, very thick, ->] (P) -- ++(0.95,-0.20);
  \node[rvector, anchor=north, inner sep=1pt] at ($(P)+(0.95,-0.22)$) {$\tvec$};
  \draw[rvector, very thick, ->] (P) -- ++(0.42,0.52);
  \node[rvector, anchor=east, inner sep=1pt] at ($(P)+(0.42,0.52)$) {$\tvecbis$};
  \draw[rtangent!85] ($(P)+(0.52,-0.11)$) arc[start angle=-12, end angle=51, radius=0.53];
  \node[rtangent!90, font=\scriptsize, anchor=west, inner sep=1pt] at ($(P)+(0.60,0.24)$) {$\phi$};

  % le couple transporte en theta_1 : meme ecart angulaire, tourne dans le
  % meme sens que la geodesique (qui vire vers le bas a droite)
  \draw[rvector, very thick, ->] (Q) -- ++(0.82,-0.51);
  \node[rvector, anchor=north, inner sep=1pt] at ($(Q)+(0.82,-0.51)$)
    {$\transport(\tvec)$};
  \draw[rvector, very thick, ->] (Q) -- ++(0.57,0.34);
  \node[rvector, anchor=south west, inner sep=1pt] at ($(Q)+(0.57,0.34)$)
    {$\transport(\tvecbis)$};
  \draw[rtangent!85] ($(Q)+(0.46,-0.29)$) arc[start angle=-32, end angle=31, radius=0.54];
  \node[rtangent!90, font=\scriptsize, anchor=west, inner sep=1pt] at ($(Q)+(0.75,-0.01)$) {$\phi$};

  \fill (P) circle (1.7pt);
  \node[anchor=north east, inner sep=2pt] at (P) {$\param_0$};
  \fill (Q) circle (1.7pt);
  \node[anchor=north west, inner sep=2pt] at (Q) {$\param_1$};
\end{tikzpicture}
